\chapter{力场、边界与可观测量}
\sourcepages{手写笔记第 5--10 页}
\section{势能模型怎样决定计算的问题}
上一章假定已知 $U(R)$。现在必须说明它包含什么物理。分子中常把键伸缩和键角变化近似成平衡附近的谐振动，例如 $U_b=k_b(r-r_0)^2/2$。它来自势能在极小值附近的 Taylor 展开，一阶导数为零，二阶导数 $k_b$ 给出刚度。这种展开适合小振幅，不能可靠描述断键。

非键作用常包含 Coulomb 项与短程排斥、色散项。以中性稀有气体的简单模型为例，Lennard--Jones 对势为
\begin{equation}
u(r)=4\epsilon[(\sigma/r)^{12}-(\sigma/r)^6],\qquad U=\sum_{i<j}u(r_{ij}).
\end{equation}
$\epsilon$ 是势阱深度，$\sigma$ 是势能过零的距离，$r_{ij}=|\bm r_i-\bm r_j|$。远距离 $r^{-6}$ 吸引反映瞬时偶极涨落之间的量子相关；诱导偶极的能量标度也提示这一幂律，但不能把它当成 London 色散的完整推导。$r^{-12}$ 是方便的经验排斥形式，不是普适定律。

令 $\bm r_{ij}=\bm r_i-\bm r_j$。利用 $\nabla_i r_{ij}=\bm r_{ij}/r_{ij}$ 求导得
\begin{equation}
\bm F_{ij}=24\epsilon[2\sigma^{12}/r_{ij}^{14}-\sigma^6/r_{ij}^8]\bm r_{ij}.
\end{equation}
当 $r<2^{1/6}\sigma$ 时力指向远离另一粒子的方向；在 $r=2^{1/6}\sigma$ 处势能取最小值 $-\epsilon$。这是检查力的符号的简单办法。

势能应在整体平移、整体转动以及相同粒子置换下不变。若机器学习模型拟合势能，再以负梯度给力，这些对称性与能量一致性仍是模型的一部分。训练误差小只说明在训练分布附近拟合好，不能代替对新构型的检验。

\section{约化单位、截断与周期边界}
Lennard--Jones 模型有自然单位 $\sigma,\epsilon,m$。令 $r^*=r/\sigma$、$T^*=\kb T/\epsilon$、$t^*=t/\tau$，其中 $\tau=\sigma\sqrt{m/\epsilon}$。代入 $m\ddot r=F$，质量和能量标度消去，得到只含无量纲变量的方程。“令 $m=\epsilon=\sigma=1$”是在选择单位，不是把实际粒子性质改成一。

直接在 $r_c$ 之外把势能设零会产生能量跳跃。平移截断 $u_s(r)=u(r)-u(r_c)$（$r<r_c$）使势连续，却仍使力跳变。若需要力也连续，可用
\begin{equation}
u_{sf}(r)=u(r)-u(r_c)-(r-r_c)u'(r_c),\quad r<r_c,
\end{equation}
外部为零。这会在截断内部也改变力，因而是模型修改。比较结果时必须报告截断方式，而不只报告截断距离。

有限盒子的表面比例过大，周期边界通过无限复制盒子模拟体相环境。立方盒边长 $L$ 时，将相对坐标逐分量变为 $\Delta x\mapsto\Delta x-L\operatorname{round}(\Delta x/L)$，得到最小镜像。短程球形截断通常要求 $r_c<L/2$，否则一个粒子的多个镜像可能进入作用范围。长程 Coulomb 相互作用不能由这一短程近似可靠处理。

计算扩散时另有一个关键区别：用于求力的坐标可以折回盒内，但位移必须用记录跨边界次数的展开坐标。如果把每次跨边界都看成跳回盒子的另一端，算出的运动距离将完全错误。

\section{怎样从快照得到结构}
平均数密度为 $\rho=N/V$。在均匀各向同性体系中，给定一个粒子后，在半径 $r$、厚度 $\dd r$ 的球壳中找到其他粒子的平均数定义为
\begin{equation}
\dd N(r)=4\pi r^2\rho g(r)\dd r.
\end{equation}
$g(r)$ 因而是相对于无关联分布的局部概率增强。计算时对每帧每个中心计数，再除以中心数、帧数、球壳体积和 $\rho$。如果只统计无序对，每个对要贡献两次。有限 $N$ 且排除中心自身时，理想气体按此归一化为 $(N-1)/N$；用 $(N-1)/V$ 可消除这一有限粒子数偏差。

第一峰表示优选近邻距离；峰后第一个谷值以内的配位数为 $4\pi\rho\int_0^{r_{\min}}r^2g(r)\dd r$。立方盒中一般只分析 $r<L/2$ 的完整球壳，避免几何截断。

\section{扩散为何与均方位移相连}
宏观上，浓度 $c$ 的梯度引起粒子流密度 $\bm j=-D\nabla c$。结合粒子数守恒 $\partial_t c+\nabla\cdot\bm j=0$，若 $D$ 为常数，得到扩散方程 $\partial_t c=D\nabla^2 c$。对起点在原点、总概率为一的传播分布 $c$，在无穷远边界项消失的条件下分部积分两次：
\begin{equation}
\frac{\dd}{\dd t}\int r^2c\,\dd^dr=D\int c\nabla^2r^2\,\dd^dr=2dD.
\end{equation}
于是三维中 $\langle|\Delta\bm r(t)|^2\rangle=6Dt$。微观轨迹的短时位移是 $\bm v(0)t$，均方位移正比于 $t^2$；只有速度记忆丧失后的扩散区才能用线性斜率估计 $D$。局域振动形成的平台也不是扩散区。

平衡体系中定义 $C_v(t)=\langle\bm v(t)\cdot\bm v(0)\rangle$。将 $\Delta\bm r=\int_0^t\bm v(s)\dd s$ 代入均方位移，并用平稳性把双重积分化成时间差积分：
\begin{equation}
\langle|\Delta\bm r(t)|^2\rangle=2\int_0^t(t-s)C_v(s)\dd s.
\end{equation}
求导并取长时间极限，即得 Green--Kubo 形式
\begin{equation}
D=\frac13\int_0^\infty C_v(s)\dd s.
\end{equation}
前提是相关函数可积且存在正常扩散极限。多粒子自扩散取各粒子自身相关的平均，不能随意加入不同粒子间的交叉项。速度相关的 Fourier 谱能显示振动频率；只有在适当谐性、质量加权和归一化条件下，才能把它直接解释为振动态密度。

\section{电输运为何不能只看单粒子扩散}
令展开坐标形成的电偶极矩为 $\bm M=\sum_iq_i\bm r_i$，总电流为 $\bm J=\dot{\bm M}$。各向同性平衡系统的线性响应给出
\begin{equation}
\sigma=\frac1{3V\kb T}\int_0^\infty\langle\bm J(t)\cdot\bm J(0)\rangle\dd t
=\lim_{t\to\infty}\frac{\dd\langle|\bm M(t)-\bm M(0)|^2\rangle/\dd t}{6V\kb T}.
\end{equation}
第二个等式由与上节相同的双重积分恒等式得到；第一个等式还依赖线性响应理论，不能仅由位移恒等式证明。展开平方会出现 $q_iq_j\langle\Delta\bm r_i\cdot\Delta\bm r_j\rangle$。忽略 $i\ne j$ 的离子协同运动，才得到 $\sigma_{\rm NE}=\sum_s\rho_s q_s^2D_s/(\kb T)$ 的 Nernst--Einstein 近似。

热导率也能由平衡涨落计算。若 $\bm J_Q$ 定义为总热流，而非单位体积流密度，则各向同性体系有
\begin{equation}
\kappa=\frac1{3V\kb T^2}\int_0^\infty\langle\bm J_Q(t)\cdot\bm J_Q(0)\rangle\dd t.
\end{equation}
具体热流必须从所用势能的能量连续性方程推导，并处理平流贡献。周期边界与多体势下，不能不加修正地把 $\sum_i e_i\bm r_i$ 的时间导数当成完整热流。

输运系数的困难往往在统计而不在公式。应改变积分上限检查平台，比较多个独立时间段，并检查盒子大小与总采样时间。相关函数的长时噪声会使积分漂移，延长积分区间并不自动提高准确度。
